%!TEX program = pdflatex
% Every content slide is GENERATED. `python presentation/make_slides.py` reads each movie's own
% spec through tools/spec_summary.py and writes slides/*.tex; this file only sets the style and
% inputs them, so nothing here has to be re-typed when a spec changes.
%
% Movies follow the deck convention: an MP4 and a same-basename PNG poster in Movies/, embedded
% as a run: link annotation that pdfpc turns into the video.
\documentclass[aspectratio=169]{beamer}

\usetheme{Janelia}
\usepackage{people}

\usepackage[english]{babel}
\usepackage[T1]{fontenc}
\usepackage[overlay,absolute]{textpos}
\setlength{\TPHorizModule}{1cm}
\setlength{\TPVertModule}{1cm}

\usepackage{amsmath,amssymb}
\usepackage{multimedia}
\usepackage{tabularx,booktabs}
\usepackage{helvet}
\usepackage{xcolor}
\usepackage{xfp}                      % \fpeval, for the divisions \fitcol needs
\usepackage{varwidth}                 % a box that takes its NATURAL width, up to a maximum

% ------------------------------------------------------------------
% BLACK SLIDES. The renders are all black-background, so the slide is black too and a frame's
% white text sits on the same field as the movie inside it -- no white margin around a black
% clip. Only the colours change; the Janelia theme's layout is untouched.
\setbeamercolor{background canvas}{bg=black}
\setbeamercolor{normal text}{fg=white,bg=black}
\setbeamercolor{structure}{fg=white}
\setbeamercolor{itemize item}{fg=white}
\usebeamercolor[fg]{normal text}
\AtBeginDocument{\color{white}}
% `gray` is redefined, not replaced: the generated slides ask for the standard colour name, so
% lightening it here is what makes a secondary line legible on black without touching the tool.
\definecolor{gray}{gray}{0.62}
% A title here is a whole sentence ("An elastic sheet pressed into a gel of 500,000 material
% points"), so the frame title is set smaller than the theme's \LARGE default: two lines of that
% do not fit the band's 1 cm and the second one fell out of it, over the slide.
\renewcommand{\frametitlesize}{\normalsize}

% ------------------------------------------------------------------
% NOTHING OVERFLOWS A SLIDE, and that is enforced rather than hoped for.
%
%   \panelbox    the height a movie or figure may occupy. A poster is bounded in BOTH
%                directions with keepaspectratio, so LaTeX shrinks whichever dimension binds
%                and a portrait render cannot grow past the frame.
%   \colheight   the height the generated spec column may occupy. `\fitcol` measures the
%                column it is given and, only if it is taller than that, scales it down to
%                exactly \colheight. Nothing is dropped: an over-full slide gets smaller type,
%                not fewer operators.
% The band is 1 cm of a 9.6 cm slide and the theme reserves nothing for it, so every generated
% frame opens with \vspace*{\bandgap}; \colheight is the rest of the slide, less the footer.
\newlength{\bandgap}
\setlength{\bandgap}{16pt}
\newlength{\panelbox}
\setlength{\panelbox}{0.72\textheight}
\newlength{\colheight}
\setlength{\colheight}{0.80\textheight}   % the slide, less the band and the footer
\newcommand{\panelboxfull}{\setlength{\panelbox}{0.78\textheight}}
\newcommand{\fitgfx}[2][\linewidth]{%
  \includegraphics[width=#1,height=\panelbox,keepaspectratio]{#2}}
% The media column carries the picture and NOTHING ELSE -- no caption, no title. What the clip
% is, is said once in the frame title; a label under it would say it twice.
% The leading \vspace{0pt} is not decoration. In a [T] column the first box is set with \vtop
% semantics, so a tall picture placed first extends UPWARD from the baseline and climbs over the
% title band; an empty first line gives the column an ordinary top to hang from.
\newcommand{\topgfx}[1]{\vtop{\hbox to \linewidth{\hss#1\hss}}}
\newcommand{\panel}[2][\linewidth]{\vspace{0pt}\topgfx{\fitgfx[#1]{#2}}}
% A movie: same fitting, over the poster PNG that pdfpc swaps for the video.
\newcommand{\playmovie}[2][\linewidth]{%
  \vspace{0pt}\topgfx{\href{run:#2.mp4?autostart&loop&noprogress}{\fitgfx[#1]{#2.png}}}}
% ONE TYPE SIZE FOR THE WHOLE DECK. `\fitfixed` is the scale every generated column is drawn
% at; 0 means "measure this slide and shrink it if it overflows", which is the fallback and
% gives a different size on every slide. presentation/fit_deck.py compiles once, reads the
% FITSCALE line each slide reports, and writes slides/fit.tex with the SMALLEST of them -- the
% densest slide's size -- so every slide is set at that one size and none has to be scaled again.
\providecommand{\fitfixed}{0}
\InputIfFileExists{slides/fit.tex}{}{}
\newsavebox{\fitbox}
\newlength{\fitht}                     % a NAMED length: beamer's columns clobber \dimen0
\makeatletter
% The scale factor is computed and applied with \scalebox rather than handing the target height
% to \resizebox{!}{...}: that form re-typesets the saved box here and keeps only its first line,
% blown up to the target -- the giant letter that was filling the right half of a slide.
% The column is measured at its NATURAL width (varwidth, not minipage) so no line ever wraps:
% "Morphogen diffusion between neighbouring cells" is one line whatever the column is worth, and
% the scale that makes it fit is taken from BOTH dimensions -- width and height, whichever binds.
\newlength{\fitwd}
\newcommand{\fitcol}[1]{%
  \sbox{\fitbox}{\begin{varwidth}{2\linewidth}#1\end{varwidth}}%
  \setlength{\fitht}{\dimexpr\ht\fitbox+\dp\fitbox\relax}%
  \setlength{\fitwd}{\wd\fitbox}%
  \edef\fit@f{\fpeval{min(1, \strip@pt\colheight / \strip@pt\fitht,
                             \strip@pt\linewidth / \strip@pt\fitwd)}}%
  \typeout{FITSCALE=\fit@f}%
  \vspace{0pt}%
  \ifdim\fitfixed pt>0pt \scalebox{\fitfixed}{\usebox{\fitbox}}%
  \else \scalebox{\fit@f}{\usebox{\fitbox}}\fi}
\makeatother
% `\srcpath` -- the file a slide was generated from, in the corner -- comes from the Janelia
% theme; it is not redefined here.
% ------------------------------------------------------------------

\title{\textbf{Plexus}}
\subtitle{A language for saying what a biological mechanism \emph{is}, precisely enough to run it}
\author{Cedric Allier}
\institute{Saalfeld lab, Janelia}
\date{October 2026}

\begin{document}

\frame{\titlepage}

% generated by presentation/make_slides.py -- the deck inputs this one file
% generated by presentation/make_slides.py from config/atlas/turing2d_gs_eta.yaml
\begin{frame}[t]{Gray-Scott chemistry on a disc of 4,000 cells}
\vspace*{\bandgap}
\begin{columns}[T,onlytextwidth]
\begin{column}{0.5\textwidth}
\playmovie{Movies/01_turing2d_gs_eta2}
\end{column}
\begin{column}{0.48\textwidth}
\fitcol{%
\vspace{6pt}
{\normalsize\textbf{entities}}\\[4pt]
{\scriptsize\begin{tabular}{@{}lr@{\hspace{7pt}}l@{}}
cell & 4,000 & \textcolor{gray}{} \\
\end{tabular}\par}
\vspace{10pt}{\normalsize\textbf{activities}}\\[4pt]
{\scriptsize\begin{tabular}{@{}c@{\hspace{5pt}}l@{}}
\raisebox{-4pt}{
\includegraphics[height=15pt]{icons/seed.png}} & \textcolor{gray}{seed} \\
 & the initial positions \\[1pt]
 & the initial morphogen \\[1pt]
\raisebox{-4pt}{\includegraphics[height=15pt]{icons/rewire.png}} & \textcolor{gray}{rewire} \\
 & neighbours within a radius \\[1pt]
\raisebox{-4pt}{
\includegraphics[height=15pt]{icons/lateral.png}} & \textcolor{gray}{lateral} \\
 & morphogen diffusion between neighbouring cells \\[1pt]
 & autocatalytic reaction \\[1pt]
\end{tabular}\par}}
\end{column}
\end{columns}
\end{frame}

% generated by presentation/make_slides.py from config/mesh_mpm/mesh_mpm_sheet_elastic.yaml
\begin{frame}[t]{An elastic sheet pressed into a gel of 500,000 material points}
\vspace*{\bandgap}
\begin{columns}[T,onlytextwidth]
\begin{column}{0.5\textwidth}
\playmovie{Movies/02_meshmpm_sheet}
\end{column}
\begin{column}{0.48\textwidth}
\fitcol{%
\vspace{6pt}
{\normalsize\textbf{entities}}\\[4pt]
{\scriptsize\begin{tabular}{@{}lr@{\hspace{7pt}}l@{}}
vertex & 4,096 & \textcolor{gray}{} \\
\hspace*{7pt}$\hookrightarrow$ half\_edge & 18,432 & \textcolor{gray}{vertex $\to$ vertex} \\
plate\_face & 2,304 & \textcolor{gray}{} \\
mpm\_particle & 500,000 & \textcolor{gray}{gel} \\
mpm\_grid & & \textcolor{gray}{field, grid 64} \\
\end{tabular}\par}
\vspace{10pt}{\normalsize\textbf{activities}}\\[4pt]
{\scriptsize\begin{tabular}{@{}c@{\hspace{5pt}}l@{}}
\raisebox{-4pt}{
\includegraphics[height=15pt]{icons/seed.png}} & \textcolor{gray}{seed} \\
 & lay out the plate \\[1pt]
\raisebox{-4pt}{
\includegraphics[height=15pt]{icons/structural.png}} & \textcolor{gray}{structural} \\
 & drive the surface at a set rate \\[1pt]
\raisebox{-4pt}{
\includegraphics[height=15pt]{icons/lateral.png}} & \textcolor{gray}{lateral} \\
 & contact with a live surface \\[1pt]
 & the material update \\[1pt]
\raisebox{-4pt}{\includegraphics[height=15pt]{icons/exchange.png}} & \textcolor{gray}{exchange} \\
 & particle to grid \\[1pt]
 & grid to particle \\[1pt]
\raisebox{-4pt}{
\includegraphics[height=15pt]{icons/field.png}} & \textcolor{gray}{field} \\
 & the grid solve \\[1pt]
\end{tabular}\par}}
\end{column}
\end{columns}
\end{frame}

% generated by presentation/make_slides.py from config/attractors/halvorsen.yaml
\begin{frame}[t]{60,000 particles riding the Halvorsen attractor}
\vspace*{\bandgap}
\begin{columns}[T,onlytextwidth]
\begin{column}{0.5\textwidth}
\playmovie{Movies/03_attractor_halvorsen}
\end{column}
\begin{column}{0.48\textwidth}
\fitcol{%
\vspace{6pt}
{\normalsize\textbf{entities}}\\[4pt]
{\scriptsize\begin{tabular}{@{}lr@{\hspace{7pt}}l@{}}
cloud & 60,000 & \textcolor{gray}{c} \\
\end{tabular}\par}
\vspace{10pt}{\normalsize\textbf{activities}}\\[4pt]
{\scriptsize\begin{tabular}{@{}c@{\hspace{5pt}}l@{}}
\raisebox{-4pt}{
\includegraphics[height=15pt]{icons/lateral.png}} & \textcolor{gray}{lateral} \\
 & flow along a strange attractor \\[1pt]
\end{tabular}\par}}
\end{column}
\end{columns}
\end{frame}

% generated by presentation/make_slides.py from config/attraction_repulsion/arbitrary_3.yaml
\begin{frame}[t]{Three particle types under one pairwise force law}
\vspace*{\bandgap}
\begin{columns}[T,onlytextwidth]
\begin{column}{0.5\textwidth}
\playmovie{Movies/04_attraction}
\end{column}
\begin{column}{0.48\textwidth}
\fitcol{%
\vspace{6pt}
{\normalsize\textbf{entities}}\\[4pt]
{\scriptsize\begin{tabular}{@{}lr@{\hspace{7pt}}l@{}}
particle & 30,000 & \textcolor{gray}{t0, t1, t2} \\
\end{tabular}\par}
\vspace{10pt}{\normalsize\textbf{activities}}\\[4pt]
{\scriptsize\begin{tabular}{@{}c@{\hspace{5pt}}l@{}}
\raisebox{-4pt}{\includegraphics[height=15pt]{icons/rewire.png}} & \textcolor{gray}{rewire} \\
 & neighbours within a radius \\[1pt]
\raisebox{-4pt}{
\includegraphics[height=15pt]{icons/lateral.png}} & \textcolor{gray}{lateral} \\
 & attraction and repulsion \\[1pt]
\end{tabular}\par}}
\end{column}
\end{columns}
\end{frame}

\input{slides/05_turing_flower_v3.tex}
% generated by presentation/make_slides.py from config/tissue/sheet_die.yaml
\begin{frame}[t]{An epithelial sheet of 400 cells extruding its dying cells}
\vspace*{\bandgap}
\begin{columns}[T,onlytextwidth]
\begin{column}{0.5\textwidth}
\playmovie{Movies/06_sheet_die_p}
\end{column}
\begin{column}{0.48\textwidth}
\fitcol{%
\vspace{6pt}
{\normalsize\textbf{entities}}\\[4pt]
{\scriptsize\begin{tabular}{@{}lr@{\hspace{7pt}}l@{}}
cell & 400 & \textcolor{gray}{of 2,000 declared} \\
vertex & 4,000 & \textcolor{gray}{} \\
\hspace*{7pt}$\hookrightarrow$ half\_edge & 16,000 & \textcolor{gray}{vertex $\to$ vertex} \\
\end{tabular}\par}
\vspace{10pt}{\normalsize\textbf{activities}}\\[4pt]
{\scriptsize\begin{tabular}{@{}c@{\hspace{5pt}}l@{}}
\raisebox{-4pt}{
\includegraphics[height=15pt]{icons/seed.png}} & \textcolor{gray}{seed} \\
 & build the cell mesh \\[1pt]
\raisebox{-4pt}{
\includegraphics[height=15pt]{icons/aggregate.png}} & \textcolor{gray}{aggregate} \\
 & cell geometry from the mesh \\[1pt]
\raisebox{-4pt}{
\includegraphics[height=15pt]{icons/lateral.png}} & \textcolor{gray}{lateral} \\
 & vertex-model shape energy \\[1pt]
\raisebox{-4pt}{
\includegraphics[height=15pt]{icons/die.png}} & \textcolor{gray}{die} \\
 & cell death \\[1pt]
\raisebox{-4pt}{\includegraphics[height=15pt]{icons/rewire.png}} & \textcolor{gray}{rewire} \\
 & T1 neighbour exchange \\[1pt]
\end{tabular}\par}}
\end{column}
\end{columns}
\end{frame}

% generated by presentation/make_slides.py from config/si_material/si_waterfall_5m.yaml
\begin{frame}[t]{A waterfall of 5,000,000 material points}
\vspace*{\bandgap}
\begin{columns}[T,onlytextwidth]
\begin{column}{0.5\textwidth}
\playmovie{Movies/07_mpm2_waterfall_5m}
\end{column}
\begin{column}{0.48\textwidth}
\fitcol{%
\vspace{6pt}
{\normalsize\textbf{entities}}\\[4pt]
{\scriptsize\begin{tabular}{@{}lr@{\hspace{7pt}}l@{}}
cell & 1 & \textcolor{gray}{water} \\
\hspace*{7pt}$\hookrightarrow$ mpm\_particle & 5,000,000 & \textcolor{gray}{5,000,000 per cell} \\
mpm\_grid & & \textcolor{gray}{field, grid 192} \\
\end{tabular}\par}
\vspace{10pt}{\normalsize\textbf{activities}}\\[4pt]
{\scriptsize\begin{tabular}{@{}c@{\hspace{5pt}}l@{}}
\raisebox{-4pt}{
\includegraphics[height=15pt]{icons/lateral.png}} & \textcolor{gray}{lateral} \\
 & gravity \\[1pt]
 & the material update \\[1pt]
 & shear viscosity \\[1pt]
\raisebox{-4pt}{\includegraphics[height=15pt]{icons/exchange.png}} & \textcolor{gray}{exchange} \\
 & particle to grid \\[1pt]
 & grid to particle \\[1pt]
\raisebox{-4pt}{
\includegraphics[height=15pt]{icons/field.png}} & \textcolor{gray}{field} \\
 & the grid solve \\[1pt]
\end{tabular}\par}}
\end{column}
\end{columns}
\end{frame}

\input{slides/08_turing2d_coupled2.tex}
\input{slides/09_meshmpm_spheroid_gel.tex}
% generated by presentation/make_slides.py from config/okuda/r020_00_ctrl.yaml
\begin{frame}[t]{A vesicle growing arms, with the purse-string tension restored}
\vspace*{\bandgap}
\begin{columns}[T,onlytextwidth]
\begin{column}{0.5\textwidth}
\playmovie{Movies/10_turing_arms_r020ctrl}
\end{column}
\begin{column}{0.48\textwidth}
\fitcol{%
\vspace{6pt}
{\normalsize\textbf{entities}}\\[4pt]
{\scriptsize\begin{tabular}{@{}lr@{\hspace{7pt}}l@{}}
cell & 2,000 & \textcolor{gray}{of 65,004 declared} \\
vertex & 130,004 & \textcolor{gray}{} \\
\hspace*{7pt}$\hookrightarrow$ half\_edge & 520,032 & \textcolor{gray}{vertex $\to$ vertex} \\
\end{tabular}\par}
\vspace{10pt}{\normalsize\textbf{activities}}\\[4pt]
{\scriptsize\begin{tabular}{@{}c@{\hspace{5pt}}l@{}}
\raisebox{-4pt}{
\includegraphics[height=15pt]{icons/seed.png}} & \textcolor{gray}{seed} \\
 & build the cell mesh \\[1pt]
 & the initial morphogen \\[1pt]
\raisebox{-4pt}{
\includegraphics[height=15pt]{icons/aggregate.png}} & \textcolor{gray}{aggregate} \\
 & cell geometry from the mesh \\[1pt]
\raisebox{-4pt}{\includegraphics[height=15pt]{icons/rewire.png}} & \textcolor{gray}{rewire} \\
 & who neighbours whom \\[1pt]
 & T1 neighbour exchange \\[1pt]
\raisebox{-4pt}{
\includegraphics[height=15pt]{icons/lateral.png}} & \textcolor{gray}{lateral} \\
 & morphogen diffusion between neighbouring cells \\[1pt]
 & autocatalytic reaction \\[1pt]
 & chemistry driven by curvature \\[1pt]
 & growth where the morphogen is high \\[1pt]
 & vertex-model shape energy \\[1pt]
 & purse-string line tension \\[1pt]
\raisebox{-4pt}{
\includegraphics[height=15pt]{icons/divide.png}} & \textcolor{gray}{divide} \\
 & cell division \\[1pt]
\end{tabular}\par}}
\end{column}
\end{columns}
\end{frame}



\end{document}
